% REVIEW DRAFT: not an input of the current published manuscript.
% This file is not proof certification; see geometry-audit-2026-11-21.json.

\section{Keeping the cutoff after zero evaluation}
\label{sec:sr-theta-simultaneous-cutoff}

The small zero evaluation must sometimes be taken in a target with one
fewer hyperbolic plane, while the first space is unchanged. A vector
in the upper full-rank cutoff need not have its zero evaluation in the
lower cutoff. The following modification, inside first-group
augmentation, supplies exactly this compatibility.

\begin{srlemma}[The finite simultaneous-cutoff modification]
\label{lem:sr-theta-simultaneous-cutoff-local}
At the selected good split place, let
\(e_0:\mathcal S_{n,d}\longrightarrow\mathcal S_{n,d-2}\)
be the actual normalized zero evaluation. Let
\(\mathcal S^o_{n,d}\) and \(\mathcal S^o_{n,d-2}\)
be the full-row-rank cutoff spaces in their specified opposite
linear models. Assume \(n\le d-2\), and that the first character
has the good-place nontrivial unitary factor \(\nu_v\).
For every \(\phi\in\mathcal S^o_{n,d}\) there is
\(\widetilde\phi\in\mathcal S^o_{n,d}\) such that
\[
 \widetilde\phi-\phi\in
       \operatorname{span}\{(\omega_v(g,1)-1)\eta:
                      \eta\in\mathcal S^o_{n,d}\},
 \qquad
              e_0\widetilde\phi\in\mathcal S^o_{n,d-2}.
 \tag{Z1}\label{eq:sr-theta-simultaneous-cutoff-local}
\]
\end{srlemma}

\begin{proof}
Use the same split Fourier identification and the same adapted
basis as in the cutoff lemma. Evaluation in the removed isotropic
\(W\) line has one restriction coordinate and one integration
coordinate. In matrix coordinates its Fourier support formula is
\[
       (e_0\phi)(C)
             =c_v\int_{k^n}\phi(0,C,z)\,dz ,
 \tag{Z2}\label{eq:sr-theta-zero-evaluation-block}
\]
with the fixed nonzero normalization \(c_v\).
The zero column and the last column are the two removed plane
coordinates; \(C\) has \(d-2\) columns.
This follows directly from invariance under \(X\otimes Ex'\):
one split idempotent acts by modulation and forces the first
column to zero, and the other acts by translation and is
integrated. The trace dual Fourier transform gives the
specified lower model. Multiplying by invertible Gram matrices
does not change either cutoff rank. The normalization of the
vacua fixes \(c_v\), whose value is irrelevant to the argument.

A matrix \((0,C,z)\) of full row rank forces
\(\operatorname{rank}C\ge n-1\).
Conversely every \(C\) of rank at least \(n-1\) admits \(z\)
for which \((0,C,z)\) has rank \(n\).
On a compact open neighborhood of such a \(C\), the same
\(z\) and a sufficiently small compact open bump about it work.
Multiply by an integral-one bump in the \(z\) variable and
a bump equal to one at the zero restriction coordinate.
A finite compact open partition proves from (Z2) that
\[
 e_0(\mathcal S^o_{n,d})
       =V:=\mathcal S\{C:\operatorname{rank}C\ge n-1\}.
\]
It also proves that \(\mathcal S^o_{n,d-2}\subset V\).

The first-group action on \(V\) is exactly the lower action:
the change of the integration variable \(z\) contributes its
Jacobian, changing the upper exponent \(d/2\) to
\((d-2)/2\). The character factor \(\nu_v\) is unchanged.
Thus \(e_0\) is an onto intertwiner
\(\mathcal S^o_{n,d}\longrightarrow V\).

Let \(V^o=\mathcal S^o_{n,d-2}\). The quotient \(V/V^o\)
is the Schwartz space on the rank-\((n-1)\) stratum.
Its first-group coinvariants vanish by the same stabilizer
calculation used in the good split coinvariant proof.
In a pivot chart the stabilizer contains a remaining
\(\mathrm{GL}_1(k)\) block with arbitrary determinant.
Its modular character and all norm factors are positive
real numbers, while the prescribed unitary factor
\(\nu_v\) is nontrivial. Choose a determinant on which it is
not \(1\); the equivariant-distribution character condition
cannot hold. Localize over the kernel Grassmannian and use
the exact Schwartz sequence of the open rank locus.
This proves
\((V/V^o)_{U(X)_v}=0\).

Right exactness of coinvariants says that
\(V=V^o+\operatorname{Aug}_{U(X)_v}V\).
Write \(e_0\phi=v^o+\sum_i(\omega_{1,v}(g_i,1)-1)v_i\).
Choose upper cutoff preimages \(\eta_i\) of \(v_i\), and put
\[
 \widetilde\phi
       =\phi-\sum_i(\omega_v(g_i,1)-1)\eta_i.
\]
The upper cutoff is stable under the first action.
Intertwining gives \(e_0\widetilde\phi=v^o\), proving (Z1).
No exactness of coinvariants on injections was used.
\end{proof}

\begin{srproposition}[An actual lower-cutoff realization of the small
\(\alpha\) component]
\label{prop:sr-theta-small-alpha-realization}
For every small polynomial upper vector in the actual cutoff,
one can change it by first-group augmentation so that its theta
period is unchanged and its zero evaluation belongs to the
actual lower cutoff. Its \(\alpha\) constant-term component is
then a fixed nonzero scalar times the ordinary convergent
lower small theta period of that zero evaluation.
In particular, after an \(A\) step the \(\alpha\) component
of the comparison space belongs to the actual lower comparison
space, including for \(q>0\).
\end{srproposition}

\begin{proof}
Present the vector as a finite sum of pure tensors with selected
finite factors in \(\mathcal S^o_{n,d}\).
Apply (Z1) at that one place to each factor. All other factors,
and all infinite polynomial Gaussians, are unchanged.
The upper cutoff is preserved, so all the modified theta
periods converge. Their differences have zero first period
by first-group invariance. Thus the actual upper theta period
is unchanged. The global zero evaluation factors into the
already chosen local evaluations, and its selected local
factor is now in the lower cutoff. Hence the entire lower
vector is in that cutoff and its lower period converges.

The full small constant-term calculation gave the combined
\(\alpha\) expression
\[
 c_0\Theta(e_0\omega(g,h)f')
      -c_1\sum_L
             \Theta(e_A\omega(\gamma_Lg,h)f').
\]
When \(e_0f'\) is in the lower cutoff, (T3) makes every
unshifted \(e_A\) term zero. The cutoff is stable under the
first action, so the same holds after each \(\gamma_Lg\).
The combined expression therefore reduces to the zero term
alone. Its integral is precisely the now convergent lower
small period of \(e_0f'\), with the fixed Haar comparison.

In an \(A\) step the large \(\alpha\) component is the actual
lower large period of its cone vector, already a polynomial
Schwartz vector. The small component just constructed is an
actual lower small period in its cutoff. Their difference is
therefore in the specified lower comparison space.
\end{proof}

\begin{srcorollary}[The boundary components vanish]
\label{cor:sr-theta-boundary-component-zero}
For every form in the actual comparison space:
if \(q=0\), its \(\beta\) line constant-term component is zero;
if \(p=0\), its \(\alpha\) line constant-term component is zero.
\end{srcorollary}

\begin{proof}
If \(q=0\), the small first space is definite, so its line
constant term has only the \(\alpha\) component.
The large \(\beta\) component is the lower zero-evaluation
period, on the unchanged first signature
\((a-p,p)\) and target signature
\((d-p-1,p-1)\).
In the canonical Fock action its first character is
\(\det^c\), \(2c=m_\alpha-m_\beta>d\).
The joint-harmonic occurrence condition proved above requires
the target negative signature to be at least \(p\).
It is \(p-1\), so this character does not occur.
Local Howe duality, or the explicit harmonic generation,
makes the first-character oscillator quotient zero.
The convergent lower period is first invariant and hence
zero. Thus the large \(\beta\) component is zero as well.

If \(p=0\), the large first space is definite and its
\(\alpha\) component is zero.
Use the preceding proposition to realize the small
\(\alpha\) component as a lower small period in its cutoff.
Its first signature is \((q,a'-q)\), its target signature is
\((d-q-1,q-1)\), and its canonical first character is
\(\det^{-c}\).
The negative-character joint-harmonic condition requires
the target negative signature to be at least \(q\).
Again it is \(q-1\), so that local first-character quotient
is zero. The lower period vanishes, proving the assertion.

These are the exact positive- and negative-character
conditions with the previously verified oscillator sign.
No inference from equal infinitesimal characters is used.
\end{proof}
